% !TEX root = ../main.tex

\section{Implementation and Performance}
\label{sec:zq-performance}

\subsection{Prototype and platform}
\label{sec:zq-platform}

Our Rust implementation uses arkworks 0.4 \cite{arkworks}, hiding KZG commitments over BN254 and a SHA-256 transcript. It maintains the committed book as orders arrive and proves both the validity of the updates and the outcome at close. To link these checks, the clearing prover reuses the final commitments from the order log. Our Noir implementation independently checks the clearing rule on a finished book identified by a salted Poseidon2 digest \cite{poseidon2}. We compile it with Nargo 1.0.0 \cite{noir} and generate zero-knowledge UltraHonk proofs with Barretenberg 6.0.0 \cite{barretenberg}. Appendices~\ref{app:zq-rust} and~\ref{app:zq-noir} describe the two implementations.

We ran both on an AMD Ryzen 3 7320U laptop on mains power, with four physical cores, eight hardware threads and 13.8 GB of memory. Our Rust implementation ran natively on Windows 11 and our Noir implementation under WSL2. We used one or eight CPU threads. The benchmarks used local test parameters. Our Rust implementation also supports importing a powers-of-tau reference string \cite{nikolaenko2024}.

\subsection{Benchmark books}
\label{sec:zq-books}

We used uniformly generated books to vary the price grid and order count independently. To test price ties on market-derived inputs, we replayed Nasdaq TotalView-ITCH order flow for AAPL \cite{itch} and selected 32 books with several volume-maximizing ticks. We also built a synthetic generator calibrated on four trading days screened for scheduled events and unusual market moves. Jointly sampling bid and ask volumes around the clearing tick produced a plateau in 4.06\% of one-second crossing books, close to the pooled replay rate of 3.79\%. Independent sampling gave 14.60\%. We drew five calibrated books at each of five sizes from 100 to 2,000 ticks.

For observations from an operating FBA, we used 362 published Cboe Europe BXE outcomes for Orange SA \cite{cboepitch,zhang2023}. Each outcome supplied an observed price and quantity for a balanced test book. Both implementations recovered those values by construction. The feed omits losing orders, so these tests cannot reconstruct Cboe's price selection. Appendix~\ref{app:zq-data} describes the replay, calibration and outcome books.

\subsection{Cost of verifying a batch}
\label{sec:zq-cost}

On five calibrated 1,000-tick books, our Rust implementation generated the clearing proof in 1.34 s on one thread and 0.37 s on eight. On the 32 replayed books, the corresponding times were 1.35 s and 0.37 s. Both workloads occupy 1,024 proof rows. A crossing batch also requires an update-validity proof, bringing the two proofs to 7,824 bytes. When no ordinary orders cross, the exchange supplies a smaller no-cross proof in place of the clearing proof.

\textbf{Table 3.} Additional proofs in our Rust implementation on uniform books with 1,000 ticks, 32-bit volumes and 250 orders. One thread. Verification includes decoding.

\begingroup
\setlength{\tabcolsep}{3pt}
\begin{longtable}{@{}lrrr@{}}
\toprule
Proof & Generation & Size & Verification \\
\midrule\endfirsthead
\toprule
Proof & Generation & Size & Verification \\
\midrule\endhead
\bottomrule\endfoot
Update validity & 0.42 s & 1,700 B & 48 ms \\
No cross & 0.40 s & 1,612 B & 12 ms \\
\end{longtable}
\endgroup

At fixed price resolution, the clearing proof keeps the same layout as the log grows. The update proof covers four logged steps per order. On eight threads, generating both proofs took 0.73 s at 1,000 orders and 1.46 s at 2,200 orders (Table 4). These figures measure proving after the close. For a 250-order batch, one thread verified the opening, log and both proofs in 79 ms from decoded inputs.

\textbf{Table 4.} Our Rust implementation on uniform 1,000-tick books, eight threads. The last column sums the clearing and update proving medians.

\begingroup
\setlength{\tabcolsep}{3pt}
\begin{longtable}{@{}lrrr@{}}
\toprule
Orders & Update rows & Update proving & Both proofs \\
\midrule\endfirsthead
\toprule
Orders & Update rows & Update proving & Both proofs \\
\midrule\endhead
\bottomrule\endfoot
250 & 1,024 & 122 ms & 0.50 s \\
1,000 & 4,096 & 347 ms & 0.73 s \\
2,200 & 16,384 & 1,083 ms & 1.46 s \\
\end{longtable}
\endgroup

\subsection{Independent clearing comparison}
\label{sec:zq-comparison}

We compare the cost of the clearing proof alone. Across the 25 calibrated and 32 replayed books, the six receipt fields exported by our Rust benchmark agreed with our Noir implementation. A separate recomputation checked all eight fields of the latter. These books contain no imbalance-only orders. The tests in Appendices~\ref{app:zq-rust} and~\ref{app:zq-noir} cover that branch.

\textbf{Table 5.} Clearing at 1,000 ticks with 32-bit volumes on one CPU thread. The code counts cover the three protocol library crates and the standalone clearing circuit, respectively.

\begingroup
\setlength{\tabcolsep}{3pt}
\begin{longtable}{@{}lrr@{}}
\toprule
Measure & Our Rust implementation & Our Noir implementation \\
\midrule\endfirsthead
\toprule
Measure & Our Rust implementation & Our Noir implementation \\
\midrule\endhead
\bottomrule\endfoot
Proof layout & 88 columns \(\times\) 1,024 rows & 100,357 gates \\
Proof generation & 1.34 s & 7.43 s \\
Proof size & 6,124 B & 14,656 B \\
Verification & 29 ms & 27 ms \\
Peak memory & about 20 MB & 157 MB \\
Source lines & 3,587 & 208 \\
\end{longtable}
\endgroup

On this workload our Rust implementation produced a proof in about one fifth of the time and used 42\% of the bytes. Verification times differed by 2 ms. We timed native function calls over 20 runs per book for our Rust implementation, including proof decoding in verification. For our Noir implementation we timed whole commands over ten runs per book. The 7.43 s sums the median witness time (1.00 s) and proof-command time (6.43 s), excluding verification-key generation. The Rust memory estimate comes from an earlier profiling run. Appendix~\ref{app:zq-bench} reports the timing distributions and scaling results.
